\chapter{Methodology}
\label{ch:methodology}

This chapter details the experiments used. I present a formal statement of the classification problem, the preprocessing pipeline used on both datasets, how symptoms are represented as features, the five classifiers I chose and why I chose them, and last but not least how I overcame problems with class imbalance, hyperparameter tuning, and evaluation metrics.

\section{Problem Formulation}
\label{sec:problem_formulation}

The task is a multi-class supervised classification task. For a patient symptom vector $\mathbf{x} \in \{0, 1\}^d$, I want to predict the disease label $y \in \{1, \ldots, C\}$ that the patient has, where $d$ denotes the number of different symptoms and $C$ the number of disease categories.

More formally, I want to learn a mapping $f: \{0, 1\}^d \rightarrow \{1, \ldots, C\}$ which minimizes the classification error rate for patients not seen before. I train the classifiers on a labeled training set $\mathcal{D} = \{(\mathbf{x}_i, y_i)\}_{i=1}^{N}$ and evaluate on a held-out test set.

This is a simplified version of the problem, in which I make three assumptions:
\begin{enumerate}
    \item The symptom representation is the same whether the symptom was self-reported or clinically observed.
    \item The datasets reasonably represent real symptom-disease patterns, though they're obviously simplified compared to actual clinical encounters.
    \item Each patient has exactly one disease label. In reality, patients often have multiple conditions simultaneously, but the datasets don't support multi-label classification.
\end{enumerate}

Figure~\ref{fig:pipeline} gives an overview of the entire experimental pipeline used in this thesis, starting from the raw data and ending with the evaluation output. Each stage is explained in detail in the remaining sections of this chapter.

\begin{figure}[H]
\centering
\begin{tikzpicture}[
    node distance=0.55cm and 0.6cm,
    every node/.style={font=\small},
    block/.style={rectangle, draw, thick, rounded corners=2pt,
        minimum width=10cm, minimum height=0.8cm, align=center, fill=blue!5},
    data/.style={rectangle, draw, thick, rounded corners=2pt,
        minimum width=10cm, minimum height=0.8cm, align=center, fill=gray!10},
    models/.style={rectangle, draw, thick, rounded corners=2pt,
        minimum width=10cm, minimum height=1.0cm, align=center, fill=orange!10},
    eval/.style={rectangle, draw, thick, rounded corners=2pt,
        minimum width=10cm, minimum height=0.8cm, align=center, fill=green!10},
    arrow/.style={-{Stealth[length=2.5mm]}, thick}
]

\node[data] (data) {%
    \textbf{Raw Data} \\
    Dataset~1: $\sim$5{,}000 records, 132 symptoms, 41 diseases \quad\textbar\quad
    Dataset~2: $\sim$246{,}000 records, 377 symptoms, 727 diseases};

\node[block, below=of data] (prep) {%
    \textbf{Preprocessing} (Sec.~\ref{sec:preprocessing}) \\
    Load CSV $\rightarrow$ deduplicate $\rightarrow$ impute missing $\rightarrow$ label encode};

\node[block, below=of prep] (split) {%
    \textbf{Stratified Train/Test Split} \quad 80\% / 20\% \quad (seed = 42)};

\node[block, below=of split] (scale) {%
    \textbf{Feature Scaling} (LR, SVM only) \quad
    \texttt{StandardScaler} fit on training set};

\node[models, below=of scale] (train) {%
    \textbf{Train 5 Classifiers} (with class weighting) \\
    Naive Bayes \quad\textbar\quad Logistic Regression \quad\textbar\quad SVM
    \quad\textbar\quad Random Forest \quad\textbar\quad XGBoost};

\node[block, below=of train] (tune) {%
    \textbf{Hyperparameter Selection} (Sec.~\ref{sec:hyperparameter_tuning}) \\
    Grid search + 5-fold stratified CV on Dataset~1 \quad\textbar\quad
    Fixed values on Dataset~2};

\node[eval, below=of tune] (eval) {%
    \textbf{Evaluation on Held-out Test Set} \\
    Accuracy \quad\textbar\quad Macro Precision \quad\textbar\quad
    Macro Recall \quad\textbar\quad Macro F1};

\node[eval, below=of eval] (out) {%
    \textbf{Outputs} \\
    Results tables \quad\textbar\quad Confusion matrices \quad\textbar\quad
    Feature importance};

\draw[arrow] (data) -- (prep);
\draw[arrow] (prep) -- (split);
\draw[arrow] (split) -- (scale);
\draw[arrow] (scale) -- (train);
\draw[arrow] (train) -- (tune);
\draw[arrow] (tune) -- (eval);
\draw[arrow] (eval) -- (out);

\end{tikzpicture}
\caption{End-to-end experimental pipeline. Each block corresponds to a stage of the methodology described in this chapter.}
\label{fig:pipeline}
\end{figure}

\section{Preprocessing Pipeline}
\label{sec:preprocessing}

To have a fair comparison across classifiers, I performed the same preprocessing pipeline on both datasets. The steps were:

\textbf{Step 1: Data Loading and Inspection.} For each data set, I used \texttt{pandas.read\_csv()} to read in the data and did initial inspection. This included identifying structural problems (such as inconsistent numbers of columns), column types (to ensure that symptoms were numeric), missing values, and duplicate records. Duplicated records may lead to over-optimistic performance if the same record is in both train and test sets, so I removed them before splitting.

\textbf{Step 2: Missing Value Handling.} In Dataset~1, I used 0 to represent the absence of a symptom, to match the binary encoding. The default is for a symptom to be considered absent if it is not specifically stated, which is a sensible default for symptom data. Dataset~2 contained few missing values (less than half a percent of all cells), and I handled them in the same manner. The alternative of dropping records with missing values would have discarded useful data.

\textbf{Step 3: Feature Encoding.} Symptoms have to be encoded as binary features. The format of Dataset~1 is already like this, with 0 or 1 in each symptom column. I use the same binary symptom table representation for Dataset~2, loading all feature columns except the disease label, replacing missing values with 0, and keeping the dataset-specific symptom names. I do not project the two data sets into a common feature space because they are compared independently.

\textbf{Step 4: Label Encoding.} I used scikit-learn's \texttt{LabelEncoder} to encode the names of the diseases as integer labels. This converts the name of every disease to its unique integer (e.g., ``diabetes'' $\rightarrow$ 0, ``pneumonia'' $\rightarrow$ 1, etc.). I saved the mapping so I could interpret the results later by converting integer predictions back to disease names. This is required because all classifiers use integer labels as their internal format.

\textbf{Step 5: Train-Test Split.} Each dataset is divided using \texttt{train\_test\_split} with \texttt{stratify=y} into an 80/20 training to testing set split. The stratify parameter ensures the class distribution in the training set is the same distribution as that in the test set. For example, if diabetes accounts for 5\% of the whole dataset, it will represent 5\% of the training data and 5\% of the test data. This is necessary for data sets that are not balanced, as a random split may accidentally include every instance of a rare disease in the training set, with no instances for testing.

I've fixed the random seed (\texttt{random\_state=42}) so that the splits are reproducible. This implies that the experiments could be rerun by anyone and the same train-test split would be obtained.

\textbf{Step 6: Feature Scaling.} For Logistic Regression and SVM, I applied \texttt{StandardScaler} to standardize features to zero mean and unit variance. These classifiers are scale-sensitive, in that if they are not scaled, then features with a larger range would dominate the model (not relevant here because all features are binary, but in the case of continuous features it is common practice).

With a binary feature, standardization has little impact because all values are already at 0 or 1, but I do put it in for consistency with standard practice. Naive Bayes and the tree-based classifiers, Random Forest and XGBoost, don't need scaling because they are scale-invariant, so I train these classifiers on the original binary features.

An important note: I fit the scaler on the training set only, then used it on both the training and test sets. Fitting the scaler on all data points prior to splitting would leak information from the test set to the training set, giving optimistic performance estimates.

\textbf{Step 7: Cross-Validation Setup.} I tuned my hyperparameters on Dataset~1 using stratified 5-fold cross-validation. Stratified splits are important because they work well with imbalanced datasets, retaining the class proportions per fold. On Dataset~2, hyperparameters are not tuned with full grid search, because full-tuning at that size would be too costly. Evaluation on both datasets is still conducted on a held-out test set.

One popular option is 5 folds, a compromise between bias and variance in the performance estimate. The fewer folds used (e.g., 3 folds), the faster the estimates will be but the noisier they become. More folds (e.g., 10) are more precise, but more expensive to compute. On the bigger Dataset~2, 5-fold cross-validation is already time-consuming, so I did not exceed it.

The entire pre-processing pipeline, software versions (Table~\ref{tab:software}) and hyperparameter search spaces (Table~\ref{tab:hyperparameters}), are documented in Appendix~\ref{app:supplementary}.

\section{Feature Representation}
\label{sec:feature_representation}

Every patient is encoded as a binary vector of dimension $d = 132$ (for Dataset~1). A patient with fever, headache and nausea would have:

\begin{equation}
\label{eq:feature_vector}
\mathbf{x} = [\underbrace{0, \ldots, 0}_{\text{other symptoms}}, \underbrace{1}_{\text{fever}}, \underbrace{0, \ldots, 0}_{\text{other symptoms}}, \underbrace{1}_{\text{headache}}, \underbrace{0, \ldots, 0}_{\text{other symptoms}}, \underbrace{1}_{\text{nausea}}, \underbrace{0, \ldots, 0}_{\text{other symptoms}}],
\end{equation}

\noindent with 1s at the location of the reported symptoms and 0s at all other locations.

This representation is similar to the bag-of-words encoding used in text classification: it captures the presence of symptoms, but not their severity, duration, or order. This is a serious drawback, but in practice binary symptom vectors are often highly successful in the classification of diseases, especially if the classifier can find useful combinations of features.

\section{Classification Models}
\label{sec:models}

For this reason I selected five classifiers, ranging from simple to complex. The selection criteria were: (1) the classifier must have an established use in medical classification, (2) the set must cover the main algorithmic paradigms (probabilistic, linear, kernel-based, and ensemble), and (3) all must have mature implementations in scikit-learn or XGBoost.

\subsection{Naive Bayes}
\label{sec:naive_bayes}

The most basic baseline is Naive Bayes. It uses Bayes' theorem assuming conditional independence of the features given the class~\citep{Bishop2006}. Since the features are binary, I am using the Bernoulli model:

\begin{equation}
\label{eq:naive_bayes}
P(y \mid \mathbf{x}) \propto P(y) \prod_{j=1}^{d} P(x_j \mid y).
\end{equation}

In reality, the independence assumption is clearly not met: the occurrences of fever and headache, for instance, are not independent. However, Naive Bayes is included as a benchmark since it trains in under a second, and gives a good lower bound against which the other more sophisticated algorithms can be compared. Implementation: \texttt{sklearn.naive\_bayes.BernoulliNB}.

\subsection{Logistic Regression}
\label{sec:logistic_regression}

Logistic Regression was selected because it is more interpretable than other models. It treats the class probabilities as a softmax over a linear combination of features~\citep{Bishop2006}:

\begin{equation}
\label{eq:logistic_regression}
P(y = k \mid \mathbf{x}) = \frac{\exp(\mathbf{w}_k^T \mathbf{x} + b_k)}{\sum_{j=1}^{C} \exp(\mathbf{w}_j^T \mathbf{x} + b_j)}.
\end{equation}

To avoid overfitting I applied L2 regularization. The benefit of this model is that the learned weights $\mathbf{w}_k$ directly represent which symptoms move the prediction toward which disease --- if a clinician asks ``why was a dermatologist recommended?'', Logistic Regression can give a definite response. Implementation: \texttt{sklearn.linear\_model.LogisticRegression} with the \texttt{lbfgs} solver.

\subsection{Support Vector Machine}
\label{sec:svm}

I added SVM as an experiment to test a margin-based classifier~\citep{Cortes1995}. For small enough datasets (such as Dataset~1) I used an RBF kernel:

\begin{equation}
\label{eq:svm_kernel}
K(\mathbf{x}_i, \mathbf{x}_j) = \exp\left(-\gamma \|\mathbf{x}_i - \mathbf{x}_j\|^2\right).
\end{equation}

In the case of Dataset~2, using a linear approximation based on stochastic gradient descent with hinge loss was much faster than a full kernel SVM would be. This maintains a practical model on the larger dataset while retaining the basic objective of the SVM. Implementation: \texttt{sklearn.svm.SVC} for Dataset~1 and \texttt{sklearn.linear\_model.SGDClassifier} for Dataset~2.

\subsection{Random Forest}
\label{sec:random_forest}

The first ensemble method that I tried is Random Forest~\citep{Breiman2001}. It creates numerous trees using random samples of data, called bootstrap samples, with each tree seeing only a random subset of features at each split:

\begin{equation}
\label{eq:random_forest}
\hat{y} = \text{mode}\{h_1(\mathbf{x}), h_2(\mathbf{x}), \ldots, h_T(\mathbf{x})\}.
\end{equation}

I selected Random Forest because it works with binary data, doesn't require any feature scaling, and gives feature importance scores (which I use later in Section~\ref{sec:feature_importance}). In preliminary tests, 200 trees yielded excellent performance, training in a few seconds. Implementation: \texttt{sklearn.ensemble.RandomForestClassifier}.

\subsection{XGBoost}
\label{sec:xgboost}

XGBoost sequentially constructs trees to correct the previous errors \citep{Chen2016}:

\begin{equation}
\label{eq:xgboost}
\mathcal{L}^{(t)} = \sum_{i=1}^{N} l(y_i, \hat{y}_i^{(t-1)} + f_t(\mathbf{x}_i)) + \Omega(f_t).
\end{equation}

I added XGBoost to see if it is worth the added complexity compared to Random Forest on symptom data. XGBoost is a very good model for tabular data, particularly in Kaggle competitions. The regularization term $\Omega$ should help on Dataset~1, which is a smaller dataset and has more risk of overfitting. Implementation: \texttt{xgboost.XGBClassifier}.

\section{Class Imbalance Handling}
\label{sec:class_imbalance}

Both datasets are imbalanced, with some diseases more prevalent than others. When not dealt with, classifiers tend to over-predict common diseases and ignore rare ones.

There were two approaches I used:

\textbf{Class Weighting:} I set \texttt{class\_weight='balanced'} for Logistic Regression, SVM and Random Forest, because each class is assigned a weight that is inversely proportional to its frequency:

\begin{equation}
\label{eq:class_weight}
w_k = \frac{N}{C \cdot n_k},
\end{equation}

\noindent Here, $N$ is the number of the entire sample, $C$ is the number of classes, and $n_k$ is the size of class $k$. This means that misclassifications on rare diseases incur a larger loss during training.

\textbf{Stratified Sampling:} Class proportions are maintained in all train-test splits and cross-validation folds.

I thought of using SMOTE~\citep{Chawla2002} but decided not to use it for this application. According to \citet{Fernandez2018}, the synthetic samples generated by SMOTE are produced by interpolation. The averaging of binary symptom vectors is problematic in the case of a binary feature space: the results of this operation are fractional symptom vectors which are not valid symptom presentations. Class weighting is a similar balancing technique that does not introduce artificial data points.

\section{Hyperparameter Tuning}
\label{sec:hyperparameter_tuning}

I optimized the hyperparameters using a grid search with stratified 5-fold cross-validation on the training set of Dataset~1. Table~\ref{tab:hyperparameters} in Appendix~\ref{app:supplementary} lists the search spaces. Dataset~2 is set with fixed parameters that were selected to ensure that the runtime is manageable.

Macro-averaged F1-score, not accuracy, was used as the scoring method for grid search. A classifier can obtain a very high accuracy on an imbalanced data set --- for instance, a disease classifier that classified all the diseases as the most common would achieve 40\% accuracy. In the Macro-F1 score, every class is given equal weight, which forces the tuning process to account for performance on rare diseases rather than only optimizing for the majority classes.

\section{Evaluation Metrics}
\label{sec:evaluation_metrics}

Each classifier was tested on four metrics, all computed on the held-out test set.

\textbf{Accuracy} is the ratio of correct predictions:

\begin{equation}
\label{eq:accuracy}
\text{Accuracy} = \frac{\text{Number of correct predictions}}{\text{Total number of predictions}}.
\end{equation}

\textbf{Precision} for each class $k$ calculates the proportion of the predicted positives which are correct:

\begin{equation}
\label{eq:precision}
\text{Precision}_k = \frac{TP_k}{TP_k + FP_k}.
\end{equation}

\textbf{Recall} is the proportion of the true positives that are correctly identified:

\begin{equation}
\label{eq:recall}
\text{Recall}_k = \frac{TP_k}{TP_k + FN_k}.
\end{equation}

\textbf{F1-score} is the harmonic mean of precision and recall:

\begin{equation}
\label{eq:f1}
\text{F1}_k = 2 \cdot \frac{\text{Precision}_k \cdot \text{Recall}_k}{\text{Precision}_k + \text{Recall}_k}.
\end{equation}

I report macro-averaged precision, recall and F1-score:

\begin{equation}
\label{eq:macro}
\text{Macro-F1} = \frac{1}{C} \sum_{k=1}^{C} \text{F1}_k.
\end{equation}

I chose macro-averaging over micro-averaging so that all disease classes are weighted equally. Under micro-averaging, the more common diseases carry most of the weight and can mask sub-optimal performance on infrequent ones. In a healthcare setting, equal weighting is more suitable, as missing a rare but serious disease is worse than misdiagnosing a common cold.

I also created confusion matrices for the best-performing classifier to find out which diseases are misclassified as each other most frequently.
